\documentclass[10pt,fleqn]{article}
\usepackage{hyperref}
\usepackage{graphicx}

\setlength{\topmargin}{-.75in}
\addtolength{\textheight}{2.00in}
\setlength{\oddsidemargin}{.00in}
\addtolength{\textwidth}{.75in}

\nofiles

\pagestyle{empty}

\setlength{\parindent}{0in}

\input{/cygdrive/m/tex/commands/commands}

\begin{document}
%%%%%%%%%%%%%%%%%%%%%%%%%%%%%%%%%%%%%%%%%%%%%%%%%%%%%%%%%%%%%%%%%%%%%%%%%%%%%%%%
%%%%%%%%%%%%%%%%%%%%%%%%%%%%%%%%%%%%%%%%%%%%%%%%%%%%%%%%%%%%%%%%%%%%%%%%%%%%%%%%
\vskip0.1in\hrule\vskip0.1in \noindent
{\bf{\Large Math 4610 Fundamentals of Numerical Analysis Tasksheet 7}}
\vskip0.1in\hrule\vskip0.1in \noindent
The problems for the Tasksheet 7 are included below. The deadline for turning
in your work on these problems will be posted on the repository. In addition,
you will turn in your work through the math4610 repository. A directory will be
constructed that will be used as a place to store your work.
\vskip0.1in\hrule\vskip0.1in \noindent
{\bf{\large Tasks}}
\vskip0.1in\hrule\vskip0.1in \noindent
\begin{trivlist}
  \item[\bf Task 1:] Write a routine to compute the solution of an upper
      triangular matrix. Test the code on the following matrix equation
         \[
            A x = b
         \]
         with
         \[
           A = (a_{i,j})
         \]
         and
         \[
           a_{i,j} = \left\lbrace
                       \begin{array}{rcl}
                         0,     & \ & j < i \\
                         \      & \ &       \\
                         i+j-1, & \ & j \geq i
                       \end{array}
                     \right.
         \]
         Finally, set \(b_i=1\).
         Hint: You should write a routine that will initialize the matrix above
         in a general way to use over and over in examples.
\vskip0.1in\hrule\vskip0.1in \noindent
  \item[\bf Task 2:] Do the same for a lower triangular system of equations.
      Create an example similar to the example in Task 1. All you need to do is
      transpose the matrix and leave the right hand side alone. Hint: It would
      be a really good idea to build code that will return the transpose of a
      matrix.
\vskip0.1in\hrule\vskip0.1in \noindent
  \item[\bf Task 3:] Write a routine that will produce a square matrix for
      testing codes in the following tasks. The input should be an integer and
      the output a square matrix of size \(n\times n\). Also, write similar,
      but separate routines that produce upper or lower triangular matrices as
      in the square matrix case. Finally, write a routine that will generate a
      diagonal matrix of a given size. Hint: These should be easy modifications
      of one another. 
\vskip0.1in\hrule\vskip0.1in \noindent
  \item[\bf Task 4:] Do the same for a diagonal systems of equations. That is
      where
         \[
           a_{i,j} = \left\lbrace
                       \begin{array}{rcl}
                         0,          & \ & j \neq i \\
                         \           & \ &       \\
                         \neq 0,     & \ & j = i \\
                       \end{array}
                     \right.
         \]
      As an example, choose random nonzero values for the diagonal entries to
      test the code. Make sure that the code is as efficient as possible.
\vskip0.1in\hrule\vskip0.1in \noindent
  \item[\bf Task 5:] Write code that will reduce a square matrix to row echelon
      form. You do not need to test for bad pivots or anything else. Write the
      code as if we know the pivots will be well behaved. Create a random matrix
\vskip0.1in\hrule\vskip0.1in \noindent
  \item[\bf Task 6:] Search the internet for sites that document the conditions
         on matrices that ensure we will be able to compute the solution of a
         linear system of equations. Write a brief paragraph (3 or 4 sentences)
         that describe your findings.  Include links to the sites you cite.
\end{trivlist}
\vskip0.1in\hrule\vskip0.1in \noindent
  \href{../../tasksheet_06/html/tasksheet_06.html}{Previous}
  \href{../../toc/md/tasksheet_toc.md}{Table of Contents} |
  \href{../../tasksheet_08/html/tasksheet_08.html}{Next}
\vskip0.1in\hrule\vskip0.1in \noindent
%%%%%%%%%%%%%%%%%%%%%%%%%%%%%%%%%%%%%%%%%%%%%%%%%%%%%%%%%%%%%%%%%%%%%%%%%%%%%%%%
%%%%%%%%%%%%%%%%%%%%%%%%%%%%%%%%%%%%%%%%%%%%%%%%%%%%%%%%%%%%%%%%%%%%%%%%%%%%%%%%
\end{document}
